\section*{Chapter \ref{ch:m3:margin-liquidation-and-loss-allocation} --- Margin, Liquidation and Loss Allocation}

\begin{solution}{exo:m3:margin-liquidation-and-loss-allocation:1}
Margin 5 per unit: liquidation at $95/0.995 = 95.48$, bankruptcy at 95.
\end{solution}

\begin{solution}{exo:m3:margin-liquidation-and-loss-allocation:2}
No: in \omterm{def:m3:margin-liquidation-and-loss-allocation:modes}{cross margin} the whole balance supports the position, so it is liquidated only when the account's total equity
falls to its maintenance margin, much further away; but a large enough fall takes the whole balance.
\end{solution}

\begin{solution}{exo:m3:margin-liquidation-and-loss-allocation:3}
The trader's margin, the \omterm{def:m3:margin-liquidation-and-loss-allocation:engine}{liquidation engine}'s sale (surplus or deficit), the \omterm{def:m3:margin-liquidation-and-loss-allocation:fund}{insurance fund}, \omterm{def:m3:margin-liquidation-and-loss-allocation:adl}{auto-deleveraging} of
profitable opposite positions, and on some venues or in the past a \omterm{def:m3:margin-liquidation-and-loss-allocation:adl}{socialised loss}.
\end{solution}

\begin{solution}{exo:m3:margin-liquidation-and-loss-allocation:4}
Linear 80.40, inverse 83.75. The inverse long's margin is in coin, which loses value as the price falls, so the equity
falls faster than the position's dollar loss alone.
\end{solution}

\begin{solution}{exo:m3:margin-liquidation-and-loss-allocation:5}
$1\%/(1 - 0.8) = 5\%$.
\end{solution}

\begin{solution}{exo:m3:margin-liquidation-and-loss-allocation:6}
In a crash its short perpetual is among the most profitable positions and can be closed by \omterm{def:m3:margin-liquidation-and-loss-allocation:adl}{auto-deleveraging}, leaving
its spot long unhedged when prices are falling fastest. It watches the ADL indicator, keeps leverage on the short low
(which lowers its rank), spreads the hedge across venues, and has a plan to re-hedge at once.
\end{solution}

\begin{solution}{exo:m3:margin-liquidation-and-loss-allocation:7}
1.65\% with leverage capped at 20 times and 5.37\% at 10 times, against 1.11\% uncapped: the cluster of high-leverage
\omterm{def:m3:margin-liquidation-and-loss-allocation:prices}{liquidation prices} near the entry is what makes a small shock dangerous, so leverage caps protect the market, not only
the traders.
\end{solution}

\begin{solution}{exo:m3:margin-liquidation-and-loss-allocation:8}
With linear permanent impact the final price is the same whatever the batch size
(\cref{prop:m3:margin-liquidation-and-loss-allocation:kappa}), and so is the \omterm{def:m3:margin-liquidation-and-loss-allocation:fund}{insurance fund}'s total, since the unit-weighted
average execution price is the same; the tutorial checks both. The answer changes only if impact is nonlinear or the book
refills between batches, and then a finer step is a different model, not a more accurate one.
\end{solution}

\begin{solution}{pb:m3:margin-liquidation-and-loss-allocation:1}
\textbf{1.} Margin 1\,000; liquidation at $24\,000/(250 \times 0.995) = 96.48$; bankruptcy at 96. \textbf{2.} $250
\times 0.5 = 125$. \textbf{3.} It receives $250 \times 0.8 = 200$. \textbf{4.} Sold above the \omterm{def:m3:margin-liquidation-and-loss-allocation:prices}{bankruptcy price}, what is left
of the trader's margin goes to the fund, which is how funds grow in normal times. \textbf{5.} The last trade can be moved
cheaply for a moment; the mark, built on a multi-venue index, cannot. \textbf{6.} 1.11\%. \textbf{7.} One: each point of
fall there triggers liquidations that cause more than a point of further fall, so the cascade runs through the zone.
\textbf{8.} 5.37\% with the cap; with the deeper book the price reaches a 10\% fall only after a 5.55\% shock, and without a jump, since $\kappa$ is then below one everywhere. \textbf{9.} With linear permanent impact the price depends only on the units sold, not on
their order. \textbf{10.} Nonlinear impact, a book that refills between batches, or engine delays that let prices gap
between batches. \textbf{11.} Near 83.83, with USD~3.21 million of deficits. \textbf{12.} About USD~1.21 million: the
deficits less the fund's USD~2 million (and a few thousand dollars of surpluses). \textbf{13.} The shorts with the highest
P\&L percentage times effective leverage: the most profitable, most leveraged winners. \textbf{14.} Its short is closed at
the \omterm{def:m3:margin-liquidation-and-loss-allocation:prices}{bankruptcy price}, so it keeps the P\&L up to that price but loses the hedge and the rest of the move, and must re-hedge
in a falling market. \textbf{15.} It would spread the loss across all winners in proportion to their gains or risk, so no
single winner loses its whole position. \textbf{16.} Caps raise the shock a market can absorb (1.11\% to 5.37\% here), at
the cost of turning away leveraged demand to other venues; a venue that caps alone may lose the flow without making the
market safer. \textbf{17.} No: a 5\% gap alone consumes it. A fund should be sized to stress scenarios of open interest,
leverage and depth, not to history. \textbf{18.} As a live risk input: reduce leverage or size on positions high in the
queue, and treat a high rank as a signal that a hedge may disappear. \textbf{19.} \emph{Named result:} a critical shock of
1.11\%; after a 5\% gap, USD~3.21 million of deficits, of which USD~2 million from the \omterm{def:m3:margin-liquidation-and-loss-allocation:fund}{insurance fund} and about USD~1.21 million
by \omterm{def:m3:margin-liquidation-and-loss-allocation:adl}{auto-deleveraging}. \textbf{20.} The loser's margin, the venue's \omterm{def:m3:margin-liquidation-and-loss-allocation:fund}{insurance fund}, and then the other winners.
\end{solution}

\begin{solution}{iq:m3:margin-liquidation-and-loss-allocation:1}
Margin $M = qE/10$. Equity $M + q(P - E)$ meets $m\,qP$ at $P = (qE - M)/(q(1 - m)) = 0.9E/(1 - m)$: about 9.5\% below entry
at $m = 0.5\%$.
\iqlookfor{equity equal to maintenance at the mark.}
\end{solution}

\begin{solution}{iq:m3:margin-liquidation-and-loss-allocation:2}
When the \omterm{def:m3:margin-liquidation-and-loss-allocation:fund}{insurance fund} cannot absorb a bankrupt position, profitable opposite positions are closed at the bankruptcy or \omterm{def:m3:perpetual-futures:index}{mark
price}. An unleveraged hedger can be deleveraged too: its hedge disappears in a crash, and its winnings are capped at a price it
did not choose.
\iqlookfor{winners bear the tail; hedges can vanish.}
\end{solution}

\begin{solution}{iq:m3:margin-liquidation-and-loss-allocation:3}
The density of \omterm{def:m3:margin-liquidation-and-loss-allocation:prices}{liquidation prices} near the current price (leverage and entry distribution), the price impact of forced sales
(depth, and whether it refills), the speed of the engines relative to arbitrage and new liquidity, and correlated positions
across venues: the local multiplier $\kappa$.
\iqlookfor{density times impact, and whether it exceeds one.}
\end{solution}

\begin{solution}{iq:m3:margin-liquidation-and-loss-allocation:4}
Liquidate early and partially (reduce to lower tiers before full liquidation), slice orders to the book's depth, use backstop
liquidity providers for large positions, trigger on a robust mark, size maintenance margins by position size and asset
liquidity, and keep the fund sized to stress scenarios.
\iqlookfor{partial, patient, and pre-arranged liquidity.}
\end{solution}

\begin{solution}{iq:m3:margin-liquidation-and-loss-allocation:5}
\omterm{def:m3:margin-liquidation-and-loss-allocation:modes}{Cross margin} shares collateral so that gains on some quotes offset losses on others and less capital is idle, but one bad
position can take the whole account. \omterm{def:m3:margin-liquidation-and-loss-allocation:modes}{Isolated margin} contains each loss but needs capital on every position and liquidates
positions a portfolio view would keep. A market maker usually runs cross or portfolio margin with its own per-instrument limits.
\iqlookfor{efficiency against containment, with internal limits.}
\end{solution}

\begin{solution}{iq:m3:margin-liquidation-and-loss-allocation:6}
From open interest changes by price level (open interest rising at a price suggests entries there), funding and long--short
ratios, leverage distributions published by venues or estimated from margin data, and liquidation feeds after the fact; combine
into a distribution of entries times a leverage distribution, and calibrate against observed cascades.
\iqlookfor{entries from open interest, leverage from any available proxy, and calibration.}
\end{solution}
